\section{Introduction}
\label{sec:intro}

If a multilingual model fails on a language because its tokenizer shreds that language into tiny pieces, we can fix the tokenizer.
If it fails because the language is structurally far from English, we cannot.
Which of the two explains zero-shot transfer loss therefore decides where practitioners should spend effort, and which languages they should flag as risky before any labeled data exists.

Multilingual encoders such as \mbert~\cite{devlin2019bert} and \xlmr~\cite{conneau2020xlmr} are routinely fine-tuned on English and applied to other languages.
The size of the resulting drop differs greatly by language~\cite{pires2019multilingual,wu2020all,hu2020xtreme}.
One appealing account points to \emph{tokenizer fertility}, the number of subword tokens per word~\cite{acs2019exploring}.
A shared vocabulary is dominated by high-resource languages, so other languages are split into more and shorter pieces, and over-segmentation hurts monolingual downstream quality~\cite{rust2021good}.
The competing account points to typological distance and pretraining corpus size~\cite{pires2019multilingual,lauscher2020zero}.

{\bf How much of the transfer drop does fertility explain?}
Existing work studies these factors separately.
What is missing is a head-to-head comparison in which fertility, typological distance and pretraining size compete to predict the \emph{same} transfer drops, across several tasks and two tokenizers, with a check of whether fertility is more than a proxy for the other two.
We test a concrete hypothesis with four parts:
(H1) the fertility ratio relative to English explains 40--60\% of the variance in transfer degradation;
(H2) it outperforms typological distance;
(H3) above 2.5$\times$ English fertility, transfer fails catastrophically (a drop of more than 30 points); and
(H4) the effect is independent of pretraining data size.

We fine-tune \mbert and \xlmr on English for natural language inference (\xnli), named entity recognition (\wikiann) and topic classification (\sib), and evaluate zero-shot on up to 71 target languages that both models saw in pretraining.
We fixed the analysis plan before training any model.
Because fertility, typology, size and script are correlated across languages, we add two designs that hold the language fixed: a within-language contrast between the two tokenizers, and an intervention that re-segments English test inputs into finer pieces.

The hypothesis is not supported as stated (see \figref{fig:r2}).
Fertility explains 0--27\% of the between-language variance in transfer drop, below the hypothesized range in all six model--task cells.
Syntactic distance from English explains 28--68\% and exceeds fertility in every cell; log pretraining size explains clearly more than fertility in four of six.
Pooled over cells, fertility reaches $R^2=0.16$ [0.07, 0.28], against 0.37 for syntactic distance and 0.32 for log size.
Once size and syntactic distance are controlled, fertility adds at most 0.05 of \rsq.
Yet fragmentation does matter in a narrower sense.
For a given language, the model whose tokenizer fragments it more loses more ($r=0.34$--$0.61$), and over-segmenting English at test time costs 10--24 points at about 1.4$\times$ fertility and 30--59 points at about 2$\times$.
Fragmentation can hurt; it is just not what separates languages that transfer well from those that do not.

In summary, our contributions are:
\begin{itemize}[leftmargin=*,itemsep=0pt,topsep=0pt]
    \item We compare fertility, \uriel typological distances and pretraining size as predictors of the same zero-shot transfer drops over 71 target languages, three tasks and two models, with bootstrap confidence intervals.
    \item We test the specific numerical claims of the fertility hypothesis and find that none holds in general: \rsq is 0--27\%, syntactic distance wins in every cell, and the 2.5$\times$ threshold is consistent with the data in one cell of six.
    \item We separate fertility from language identity with a within-language tokenizer contrast and an over-segmentation intervention, and show that fragmentation has a real cost when it is unfamiliar to the model.
    \item We report an exploratory lead: a tokens-per-parallel-sentence measure reaches $R^2=0.42$--$0.51$ on \sib and survives controls there, but not on NER or \xnli.
\end{itemize}

We review related work in \secref{sec:related}, describe the setup in \secref{sec:method}, present results in \secref{sec:results}, and discuss interpretation and limitations in \secref{sec:discussion}.
